%% kapa: exam=1 round=31 year=2020 period=1 form=A subject=civil_law subject_ko=민법 seq=7 number=7
%% source: Q-Net 감정평가사 기출문제 (cst00309) artlSeq=5207219 file=extracted/1008/5207219/A형/A형/2020년도 제31회 감정평가사 1차 시험_1교시_A형.pdf
%% conversion: kordoc 4.24.51 -> tools/gen1.py (auto); review=proofread-sonnet; flags=none
\begin{question}{7}
물건에 관한 설명으로 옳지 않은 것은? (다툼이 있으면 판례에 따름)
\begin{choices}
  \item 종물은 주물소유자의 상용에 공여된 물건을 말한다.
  \item 주물과 다른 사람의 소유에 속하는 물건은 종물이 될 수 없다.
  \item 주물과 종물의 관계에 관한 법리는 권리 상호간에도 적용된다.
  \item 저당권의 효력이 종물에 미친다는 규정은 종물은 주물의 처분에 따른다는 것과 이론적 기초를 같이 한다.
  \item 토지의 개수는 지적공부의 등록단위가 되는 필(筆)을 표준으로 한다.
\end{choices}
\end{question}
\answer{1}
